% <Course> <TF session / notes>: <topic>. A reading handout.
% Design: tfhandout.sty (~/Library/texmf/tex/latex/tfhandout/).
% Compile with XeLaTeX: latexmk -xelatex handout.tex
\documentclass{article}
\usepackage{tfhandout}
\hypersetup{pdftitle={<Course> <Session>: <Topic>}}

\title{\textsc{<TF Session N>} \\ <Plain Economist-Style Title>\thanks{Based on the handouts from Prof. <Name> and my notes from econometrics.}}
\author{TF: Rony Rodríguez-Ramírez}
\date{<Course> $\cdot$ <School>}

\begin{document}
\maketitle

\section{<First section: start here, no summary paragraph>}\label{sec:first}
Let's assume that \ldots{} Therefore, \ldots
\begin{empheq}[box=\tcbhighmath]{equation}\label{eq:key}
y = \beta_0 + \beta_1 x + u
\end{empheq}

\begin{defin}[<Term>, <source>]
<Definition.>
\end{defin}

\begin{keybox}[title={A student's question}]
\emph{``<Question, verbatim.>''}\par\smallskip
My short answer: \ldots
\end{keybox}

% Tables: an R script writes a complete tabular fragment to Tabs/; never edit it by hand.
\begin{table}[tbp]
  \centering
  \singlespacing
  \adjustbox{max width=\textwidth}{%
    \begin{threeparttable}
      \caption{<Caption on top>}\label{tab:main}
      \input{Tabs/main.tex}
      \begin{tablenotes}
        \setlength\labelsep{0pt}
        \footnotesize
        \item \textit{Notes}: <Sample, specification, standard errors.>
      \end{tablenotes}
    \end{threeparttable}
  }
\end{table}

% Figures: caption on top, plain-language notes below.
\begin{figure}[tbp]
\centering
\caption{<Caption on top>}\label{fig:main}
% \includegraphics[width=0.8\textwidth]{Figs/<file>.pdf}  or TikZ / pgfplots
\fnote{\textit{Notes}: <How to read the figure.>}
\end{figure}

% \bibliography{refs}   % only if something is cited

% \newpage
% \appendix
% \section{R code}\label{app:code}
% \lstinputlisting[style=rcode]{../code/<script>.R}
\end{document}
